\chapter{Uncertainty}

\ex{
Consider an investor who
must decide how much of their initial wealth \(w\) to put into a risky
asset. The risky asset can have any of the positive or negative rates of
return \(r_i\) with probabilities \(p_i\), \(i = 1, 2,\cdots, n\).
Suppose the investor is an expected utility maximizer and their utility
for \(x\) amount of money for sure can be represented by a twice
continuously differentiable, strictly increasing and strictly concave
utility function \(u\pr{x}\). Let \(a^*\) be the investor's optimal
amount of money to put in the risky asset. Give a necessary and
sufficient condition for the investor to have strict incentives to invest
in the risky asset, that is, \(a^*>0\) and $a^*$ is strictly preferred to
\(a = 0\).
}

\sol{
Since the investor is an expected utility maximizer, optimization problem is given by:
\[\max_{0\le a\le w} U\pr{a} = \sum_{i = 1}^n p_i\cdot u\pr{w+ar_i}\]

A sufficient condition for $a^*>0$ is
\(U'\pr{0} = u'\pr{w}\cdot \sum_{i = 1}^n p_ir_i>0\), that is, the
expected rate of return \(\sum_{i = 1}^n p_ir_i>0\).

Next suppose \(\sum_{i  =1}^np_ir_i\le 0\). Since \(u\pr{\cdot}\) is
strictly concave,
\(U''\pr{a} = \sum_{i = 1}^n p_i\cdot r_i^2 \cdot u''\pr{w+ar_i}\le 0\),
which implies \(U\pr{\cdot}\) is also strictly concave. Consequently,
\(U'\pr{a}\le U'\pr{0}\le 0\) for $a\ge 0$. It follows that the condition is also
necessary.
}

Convenient as the expected utility representation is, it is rather ad-hoc.
Similar to the essence of previous works, we want to understand what
restrictions on the agent's preference relation induce an expected
utility representation.

\hypertarget{basic-setups}{%
\section{Setups}\label{basic-setups}}

\defn{Simple Lottery; Compound Lottery}{
Let \(X = \brace{\x_1,\x_2,\cdots,\x_n}\) denote a
finite set of \textit{sure} outcomes (i.e., without uncertainty).
\begin{itemize}
\item A \textit{simple lottery}
  \(\p = p_1\circ \x_1 + p_2 \circ \x_2 + \cdots + p_n\circ\x_n\) ($p_1+p_2+\cdots+p_n = 1$) is a
  probability distribution over a finite number of sure outcomes $\brace{\x_1,\x_2,\cdots,\x_n}$. When there is no
  confusion, we will also use \(\p = \brace{\seq[p]}\) to denote such a
  simple lottery.
\item A \textit{compound lottery} \(\sum_{j = 1}^k\alpha_j\p^j\) ($\alpha_1+\alpha_2+\cdots+\alpha_k = 1$) is a probability distribution over a finite number
  of simple lotteries $\brace{\p^1,\p^2,\cdots,\p^k}$.
\end{itemize}
}

To facilitate our analysis, we introduce two assumptions with regard to the way in which utility maximizer regard the lotteries.

\asm{}{
\begin{enumerate}

\item
  The probabilities are \textbf{objective}: the probabilities are regarded as
  objective or exogenously given facts by the decision-maker (in contrast to subjective
  evaluation).
\item
  The decision-maker is a \textbf{consequentialist}: decision-makers only care
  about the outcomes, instead of how lotteries are mixed.

  Under this assumption, the compound lottery
  \(\sum_{j = 1}^k \alpha_j\p^j\) is identical to the simple lottery it
  induces, i.e.,
  \[\sum_{j = 1}^k\alpha_j\p^j \iff \pr{\sum_{j = 1}^k\alpha_j\p_1^j,\sum_{j = 1}^k\alpha_j\p_2^j,\cdots,\sum_{j = 1}^k\alpha_j\p_n^j}\]
\end{enumerate}
}

\rmk{
 The first assumption can be relaxed.
}

Given the two assumptions, we can restrict our focus on the space of simple
lotteries \(\mathcal P = \Delta\pr{X}\).

\defn{Space of Simple Lotteries}{
The space of simple lotteries, $\Delta\pr{x}$, is defined as
\[\Delta\pr{X} = \brace{\pr{p_1,\cdots ,p_n}:p_i\ge 0\ \ \forall i,p_1+\cdots + p_n = 1}.\]
}

\noindent Consider the agent's preference relation \(\pref\) over
\(\P\). To ensure a utility representation \(U:\P\to \RR\), we maintain
the three basic assumptions as in the certainty benchmark.

\begin{itemize}
\item
  Completeness: For any \(\p^1,\p^2\in \P\), either \(\p^1\pref\p^2\) or
  \(\p^2\pref\p^1\) (or both).
\item
  Transitivity: For any \(\p^1,\p^2,\p^3\in\P\), if \(\p^1\pref\p^2\)
  and \(\p^2\pref\p^3\), then \(\p^1\pref\p^3\).
\item
  Continuity: For any two sequences
  \(\pr{\p^i}_{i = 1}^n,\pr{\q^i}_{i = 1}^n\in \P\), if
  \(\p^i\pref\q^i\pr{\forall i = 1,2,\cdots,n}\) and
  \(\lim_{n\to\infty}\p^i = \p^*\), \(\lim_{n\to\infty}\q^i = \q^*\),
  then \(\p^*\pref\q^*\).
\end{itemize}

\rmkb{
There is an alternative definition of continuity.
Let \(\pref\) be a complete, transitive preference
relation on \(\P\). $\pref$ is \textit{continuous} if for any \(\p^1,\p^2,\p^3\in \P\) such that
\(\p^1\pref\p^2\pref\p^3\), there exists \(t\in \br{0,1}\) such that
\(\p^2\sim t\p^1+\pr{1-t}\p^3\).
}

To ensure an expected utility representation, we need an additional
restriction.

\defn{Independence}{
A preference relation \(\pref\) on
\(\P\) satisfies \emph{independence} if for any \(\p^1,\p^2,\p^3\in \P\)
and \(t\in\pr{0,1}\),
\[\p^1\pref \p^2\iff t\p^1+\pr{1-t}\p^3\pref t\p^2+\pr{1-t}\p^3\]
}

Intuitively, this assumption requires that the relative ranking of two
lotteries will not be affected by the mixture with a third lottery.
Notice that both \(t\p^1+\pr{1-t}\p^3\) and \(t\p^2+\pr{1-t}\p^3\) are
\emph{compound} lotteries, which should be understood as the
simple lotteries they induce.

Moreover, if the preference relation \(\pref\) on \(\P\) satisfies
independence, then for any \(\p^1,\p^2,\p^3\in \P\) and
\(t\in \pr{0,1}\), we have

\begin{itemize}
\item
  \(\p^1\succ \p^2\iff t\p^1+\pr{1-t}\p^3\succ t\p^2+\pr{1-t}\p^3\)
\item
  \(\p^1\sim \p^2\iff t\p^1+\pr{1-t}\p^3\sim t\p^2+\pr{1-t}\p^3\)
\item
  If \(\p^1\pref\p^2\) and \(\p^3\pref \p^4\), then
  \(t\p^1+\pr{1-t}\p^3\pref t\p^2+\pr{1-t}\p^4\).
\end{itemize}

In summary, in this chapter of uncertainty, we need the following four assumptions to guarantee expected utility representation.

\asm{}{
\begin{itemize}
\item
  \textbf{Completeness}: For any \(\p^1,\p^2\in \P\), either \(\p^1\pref\p^2\) or
  \(\p^2\pref\p^1\) (or both).
\item
  \textbf{Transitivity}: For any \(\p^1,\p^2,\p^3\in\P\), if \(\p^1\pref\p^2\)
  and \(\p^2\pref\p^3\), then \(\p^1\pref\p^3\).
\item
  \textbf{Continuity}: For any two sequences
  \(\pr{\p^i}_{i = 1}^n,\pr{\q^i}_{i = 1}^n\in \P\), if
  \(\p^i\pref\q^i\pr{\forall i = 1,2,\cdots,n}\) and
  \(\lim_{n\to\infty}\p^i = \p^*\), \(\lim_{n\to\infty}\q^i = \q^*\),
  then \(\p^*\pref\q^*\).
  \item \textbf{Independence}: For any \(\p^1,\p^2,\p^3\in \P\) such that
\(\p^1\pref\p^2\pref\p^3\), there exists \(t\in \br{0,1}\) such that
\(\p^2\sim t\p^1+\pr{1-t}\p^3\).
\end{itemize}
}

\hypertarget{expected-utility-representation}{%
\section{Expected Utility
Representation}\label{expected-utility-representation}}

Before starting the representation theorem, we first formalize what we
mean by an expected utility function.

\defn{Expected Utility Form}{
A utility function \(U:\P\to\RR\) has an
\emph{expected utility form} (or a \textit{von Neumann-Morgenstern expected
utility function}) if there is an assignment of numbers
(\(u_1,u_2,\cdots,u_n\)) to each of the \(n\) sure outcomes
\(\x_1,\x_2,\cdots,\x_n\) such that for any \(\p\in \P\),
\[U\pr{\p} = \sum_{i = 1}^n p_iu_i.\]
}

A sure outcome \(\x_i\) can be represented by a simple lottery \(\textbf e_i = \pr{0,\cdots ,0,1,0,\cdots, 0}\) and we have
\(U\pr{\textbf e_i} = u(\x_i) = u_i\). Notice that an expected utility function can be written as
\[U\pr{\p} = \sum_{i = 1}^n p_iU\pr{\textbf e_i}.\]

Thus, $U\pr{\p}$ is linear in the probabilities. The following proposition reveals that this observation rings true more generally.

\propp{}{
A linear function \(U:\P\to \RR\) has an
expected utility form if and only if it is \emph{linear}, that is, for
any \(k\) (simple) lotteries \(\p^j\in \P,\pr{j = 1,2,\cdots,k}\) and
associated probabilities \(\alpha_j\ge 0\) with
\(\sum_{j = 1}^k\alpha_j = 1\), we have
\[U\pr{\sum_{j = 1}^k \alpha_j\p^j} = \sum_{j = 1}^k\alpha_jU\pr{\p^j}.\]
}{
\begin{itemize}
\item
  ``If" part

  Let $\p^j = \e_j$ for all $j = 1,2,\cdots,k$. Thus, $\p = \sum_{j = 1}^k \alpha_j\p^j = \pr{\alpha_1,\alpha_2,\cdots,\alpha_k}$. The condition of $U$ being linear is then transformed to
    \[U\pr{\p} = \sum_{j = 1}^k p_jU\pr{\textbf e_j} = \sum_{j = 1}^k\alpha_ju_j.\]
    This directly implies that $U$ has an expected utility form.
\item
  ``Only if" part

  \begin{itemize}
  \item
    Since the decision-maker is a consequentialist, the compound lottery
    \(\sum_{j =1}^k \alpha_j\circ \p^j\) is identical to the simple
    lottery it induces,
    \[\alpha_1\circ \p^1+\alpha_2\circ \p^2+\cdots+\alpha_k\circ \p^k
    \iff
    \left(\sum_{j = 1}^k \alpha_jp_1^j,\sum_{j = 1}^k \alpha_jp_2^j, \cdots,\sum_{j = 1}^k \alpha_jp_n^j\right)\]
  \item
    Since the utility function has the expected utility form, we have
    \[U\pr{\sum_{j = 1}^k \alpha_jp_1^j,\sum_{j = 1}^k \alpha_jp_2^j, \cdots,\sum_{j = 1}^k \alpha_jp_n^j} = \sum_{i = 1}^n\pr{\sum_{j = 1}^k\alpha_jp^j_i}u_i\]
  \item
    On the other hand, \(U\pr{\p^j} = \sum_{i = 1}^n p_iu_i\), so
    \[\sum_{j = 1}^k \alpha_j U\pr{\p^j} = \sum_{j = 1}^k \alpha_j \pr{\sum_{i = 1}^n p_i^ju_i}\]
  \item
    Jointly, we have proved the result:
    \[\ba
    U\pr{\sum_{j = 1}^k\alpha_j\p^j} & = 
    U\pr{\sum_{j = 1}^k \alpha_jp_1^j,\sum_{j = 1}^k \alpha_jp_2^j, \cdots,\sum_{j = 1}^k \alpha_jp_n^j}\\
    & = \sum_{i = 1}^n\pr{\sum_{j = 1}^k\alpha_jp^j_i}u_i \\
    & = \sum_{j = 1}^k \alpha_j U\pr{\p^j}
    \ea\]
  \end{itemize}
\end{itemize}
}

\thmp{}{
A complete and transitive preference relation
\(\pref\) on \(\P\) satisfies continuity and independence if and only if
it admits an expected utility representation \(U:\P\to \RR\).
}{
The ``if" part is easy to verify. We mainly focus on the "only if" part.

Suppose the complete and transitive preference relation $\pref$ satisfies continuity and
    independence. We construct an expected utility function
    \(U\pr{\p} = \sum_{i = 1}^n p_iu_i\) that represents \(\pref\) as
    follows:

    \begin{itemize}
    \item
      First, with a slight abuse of notation, interpret the consequences
      \(x_1,\cdots,x_n\) as degenerate lotteries, so that \(x_i\) puts
      probability \(1\) on \(x_i\) and probability \(0\) on all other
      consequences. Then any lottery \(p\) can be viewed as a mixture of
      those degenerate lotteries: \(p = \sum_{i}p_ix_i\).
    \item
      Label the consequences so that (when viewed as degenerate
      lotteries) \(x_1\) is the best and \(x_n\) is the worst, i.e.,
      \(x_1\pref x_i\pref x_n\) for all \(i\).
    \item
      By continuity, for each consequence \(x_i\), there is some
      \(\l_i\in \br{0,1}\) such that \(x_i\sim \l_ix_1+\pr{1-\l_i}x_n\).
    \item
      By independence (iterated \(n\) times for each consequence), for
      any lottery \(p = \sum_{i}p_ix_i\),
      \[p\sim \sum_i p_i\pr{\l_ix_1+\pr{1-\l_i}x_n} = \pr{\sum_ip_i\l_i}x_1+\pr{1-\sum_i p_i\l_i}x_n\]
      \begin{itemize}
      \item
        If \(x_1\sim x_n\) then by independence, for all \(p\) we have
        \[p\sim \pr{\sum_i p_i\l_i}x_1+\pr{1-\sum_ip_i\l_i}x_1 = x_1\]
        Hence, letting \(u_i = 0\) for all \(i\) gives an expected utility representation of \(\pref\).
      \item
        If \(x_1\succ x_n\), then for \(1\ge \l\ge \l'\ge 0\), letting
        \(\delta = \l-\l'\in\br{0,1}\) and
        \(\hat p = \frac{\l'}{1-\delta} x_1+\frac{1-\l}{1-\delta}x_n\)
        and using the Independence Axiom,
        \[\ba
        \l x_1+\pr{1-\l}x_n & = \delta x_1 + \pr{1-\delta}\hat p\\
        & \succ \delta x_1 + \pr{1-\delta }\hat p\\
        & = \l'x_n+\pr{1-\l'}x_n
        \ea\]
        Thus, letting \(u_i = \l_i\) yields an expected utility representation of
        \(\pref\).
      \end{itemize}
    \end{itemize}
}

\hypertarget{measures-of-risk}{%
\section{Measures of Risk}\label{measures-of-risk}}

\hypertarget{risk-attitudes}{%
\subsection{Risk Attitudes}\label{risk-attitudes}}

Let the set of sure outcomes \(X = \RR\). For any sure outcome $x\in X$, its utility is determined by utility function $u:X\to \RR$. A lottery is fully characterized by a corresponding distribution function \(F\), whose expected utility is then 
\(U\pr{F} = \int_Xu\pr{x}\text d F\pr{x}\). For simplicity, suppose
\(u\pr{\cdot}\) is strictly increasing and continuous and
\(U\pr{F} = \E[F]{u\pr{x}}<+\infty\).

\rmk{
Here we naturally extend the notion of lottery to $X = \RR$, the infinite set, and avoid any technical issues related to the definition of lotteries on finite-many sure outcomes.
}

\defn{Risk Averse}{
For any non-degenerate lottery $F$, define its average payoff $\delta_{\text{E}_F}$ as
\[\delta_{\text{E}_F} = \int_{X}x\text dF\pr x.\]
A decision-maker is strictly \emph{risk-averse} if for any
non-degenerate lottery \(F\), the sure outcome $\delta_{\text{E}_F}$ is strictly preferred to the lottery $F$, i.e., \(\delta_{\text{E}_F}\succ F\). Or in utility terms, 
\[u\pr{\delta_{\text E_F}} = u\pr{\int_{X}x\text dF\pr x} > U\pr{F} = \int_{X}u\pr{x}\text dF\pr{x}\]
}

\cor{}{
A decision-maker is (strictly) risk-averse if and only
if \(u\pr{\cdot}\) is (strictly) concave.
}

Intuitively, the marginal utility of an 
additional unit of money is decreasing for a risk-averse individual.

\textit{Risk-loving} and \textit{risk-neutral} attitudes can
be similarly defined.

We have defined risk attitudes based on how agents evaluate the utility of averaged payoff against the averaged utility. Apart from determining risk attitudes, we develop a cursory measure, certainty equivalent, to capture such ``deviation".

\defn{Certainty Equivalent}{
The \emph{certainty
equivalent} \(c\pr{F,u}\) of a money lottery is defined as
\[u\pr{c\pr{F,u}} = U\pr{F}.\]
}

Intuitively, \(c\pr{F,u}\) is the sure outcome perceived as equivalent to the risk lottery.

\rmk{
Certainty equivalent depends on initial wealth. For
simplicity, we assume \(w_0 = 0\), i.e., the agent does not have initial
wealth; then \(c\pr{F,u}\le \E{F}\) if and only if the agent is
risk-averse.
}

\hypertarget{measure-of-risk}{%
\subsection{Measures of Risk Aversion}\label{measure-of-risk}}

We know that an agent is (strictly) risk-averse if and only if
\(u\pr{\cdot}\) is (strictly) concave. Recall the cardinal
interpretation of vNM expected utility. Can we use \(u\pr{\cdot}\) to
quantify and then compare different levels of risk aversion? Therefore, we introduce the
measures of risk aversion.\\

\noindent Consider a risk-averse agent who has an initial wealth \(w\) and faces a
(small) fair gamble.

\begin{itemize}
\item
  Scenario 1: The gamble \(\tilde\ve\) is measured in terms of monetary unit.
  Then consider how large $a$ has to be to make
  \[u\pr{w-a} = U\pr{w+\tilde \ve}.\]
\item
  Scenario 2: The gamble \(\tilde\delta\) is measured in terms of  percentage of
  the initial wealth. Then consider how large $r$ has to be to ensure
  \[u\pr{\pr{1-r}w} = U\pr{\pr{1+\tilde\delta}w}.\]
\end{itemize}

Since the agent is risk-averse, \(a,r>0\). Intuitively, \(a\) measures
how much money the agent is willing to give up to avoid the gamble;
\(r\) measures the percentage of initial wealth the agent is willing to
give up to avoid the gamble. If the agent were risk-neutral,
\(a = r= 0\). As a result, given initial wealth \(w\) and the small fair gamble, both \(a\) and \(r\)
measure the agent's level of risk aversion.

\paragraph{Coefficient of Absolute Risk Aversion}

When \(\tilde\ve\) is small, \(a\) is small. By Taylor expansion and definition of expected utility:
\[\ba
u\pr{w-a}&\approx u\pr{w}-u'\pr{w}\cdot a\\
U\pr{w+\tilde\ve} & = \E[\ve]{u\pr{w+\ve}} = \int_{\ve}u\pr{w+\ve}\text d\ve\\
& \approx \int_{\ve} \br{u\pr{w}+u'\pr{w}\cdot\ve+\frac12 u''\pr{w}\cdot\ve^2} \text d\ve\\
& = u\pr{w} \int_{\ve} 1 \text d\ve + u'\pr{w} \int_{\ve} \ve \text d\ve +\frac12  u''\pr{w}\int_{\ve} \ve^2 \text d\ve\\
& = u\pr{w} + \frac12 u''\pr{w}\cdot \var{\tilde\ve}
\ea\]

It follows that
\[\ba
u\pr{w}-u'\pr{w}\cdot a & = u\pr{w} + \frac12 u''\pr{w}\cdot \var{\tilde\ve}\\
\too a & \approx -\frac12\cdot \frac{u''\pr{w}}{u'\pr{w}}\cdot \var{\tilde\ve}
\ea\]

\defn{Coefficient of Absolute Risk Aversion}{
Suppose the agent has initial wealth $x>0$, and the utility function $u:\RR\to \RR$ is twice continuously differentiable. Then the  \textit{coefficient of absolute risk aversion}, $A\pr{x,u}$, is defined as
\[A\pr{x,u} = -\dfrac{u''\pr{x}}{u'\pr{x}}.\]
}

\paragraph{Coefficient of Relative Risk Aversion}

Similarly, we use Taylor expansion and definition of expected utility to compute
\(r\):
\[\ba
u\pr{w\pr{1-r}}& \approx u\pr{w} - u'\pr{w}\cdot \pr{wr}\\
U\pr{w+w\tilde\delta} & = \E[G]{u\pr{w+w\delta}} = \int_{\delta}u\pr{w+w\delta}\text dG\pr{\delta}\\
& \approx \int \br{u\pr{w}+u'\pr{w}\cdot\pr{w\delta}+\frac12 u''\pr{w}\cdot\pr{w\delta}^2} \text d G\pr{\delta}\\
& = u\pr{w} \int 1 \text dG\pr{\delta} + u'\pr{w} \int_{\delta} w\delta \text dG\pr{\delta} +\frac12  u''\pr{w}\int_{\delta} \pr{w\delta}^2 \text dG\pr{\delta}\\
& = u\pr{w} + \frac12 u''\pr{w}\cdot w^2\var{\tilde\delta}
\ea\]

It follows that
\[\ba
u\pr{\pr{1-r}w} = U\pr{\pr{1+\tilde\delta}w}\\
\too r \approx -\frac{wu''\pr{w}}{2u'\pr{w}}\var{\tilde\delta}
\ea\]

\defn{Coefficient of Relative Risk Aversion}{
Suppose the agent has initial wealth $x>0$, and the utility function $u:\RR\to \RR$ is twice continuously differentiable. Then the  \textit{coefficient of relative risk aversion}, $R\pr{x,u}$, is defined as
\[A\pr{x,u} = -\dfrac{xu''\pr{x}}{u'\pr{x}}.\]
}

\propp{Equivalence of Different Measures of Risk Aversion}{
Suppose utility functions $u$ and $v$ are strictly increasing and twice differentiable. The following definitions of an agent characterized by \(u\) being more risk averse than another agent characterized by \(v\) are equivalent:
\begin{enumerate}
\item
  For any lottery \(F\) and sure outcome \(\delta_X\). if
  \(F\pref_u\delta_X\), then \(F\pref_v \delta_X\).
\item
  For any lottery \(F\), \(c\pr{F,u}\le c\pr{F,v}\).
\item
  The function \(u\) is ``more concave" than \(v\), that is, there exists
  some increasing and concave function \(g\) such that \(u = g\circ v\).
\item
  \(r\pr{x} = \dfrac{u'\pr{x}}{v'\pr{x}}\) is non-increasing in \(x\).
\item
  For any \(x\),
  \(A\pr{x,u} = -\dfrac{u''\pr{x}}{u'\pr{x}}\ge -\dfrac{v''\pr{x}}{v'\pr{x}} = A\pr{x,v}\).
\end{enumerate}
}{
\begin{itemize}
\item
  (1)\(\iff\)(2)
  \begin{itemize}
  \item
    (1) \(\too\) (2): By definition of certainty equivalent,
    \(\delta_{c\pr{F,u}}\sim_u F\). By condition (1), we have
    \(\delta_{c\pr{F,v}}\sim_v F\pref_v \delta_{c\pr{F,u}}\). Recall
    that \(v\pr{\cdot}\) is strictly increasing, so we have
    \(c\pr{F,u}\le c\pr{F,v}\).
  \item
    (2) \(\too\) (1): If \(\delta_{c\pr{F,u}}\sim F\pref_u \delta_x\),
    then by strictly-increasing property of \(u\), it follows that
    \(c\pr{F,u}\ge x\). So
    \(c\pr{F,v}\ge c\pr{F,u}\ge x\too c\pr{F,v}\ge x\). Again by $v$ being strictly increasing, we have
    \(c\pr{F,v}\sim F\pref_v \delta_x\).
  \end{itemize}
\item
  (2)\(\iff\)(3)
  \begin{itemize}
  \item
    Since \(u\) and \(v\) are strictly increasing, \(v^{-1}\) is well
    defined and \(g = u\circ v^{-1}\) is strictly increasing.
  \item
    Again by strict monotonicity of \(u\), we have
    \(u\pr{c\pr{F,u}}\le u\pr{c\pr{F,v}}\) for any lottery \(F\).
    \[\ba
    u\pr{c\pr{F,u}} & = U(F) = \int _X u\pr{x}\text dF\pr{x} = \int_X g\pr{v\pr{x}}\text dF\pr{x}\\
    u\pr{c\pr{F,v}} & = g\pr{v\pr{c\pr{F,v}}} = g\pr{\int_X v\pr{x}\text dF\pr{x}}
    \ea\]
  \item
    By Jensen's inequality,
    \(\int_X g\pr{v\pr{x}}\text dF\pr{x}\le g\pr{\int_X v\pr{x}\text dF\pr{x}}\)
    if and only if \(g\) is concave.
  \end{itemize}
\item
  (3)\(\iff\)(4)

  \begin{itemize}
  \item
    Since \(u\) and \(v\) are strictly increasing, \(v^{-1}\) is well
    defined and \(g = u\circ v^{-1}\) is strictly increasing.
  \item
    When \(u\) and \(v\) are both differentiable,
    \(u'\pr{x} = g'\pr{v\pr{x}}\cdot v'\pr{x}\too \frac{u'\pr{x}}{v'\pr{x}} = g'\pr{v\pr{x}}\).
  \item
    Since \(v\) is strictly increasing, \(\frac{u'\pr{x}}{v'\pr{x}}\) is
    non-increasing in \(x\) if and only if \(g'\) is non-increasing, that is, \(g\)
    is concave.
  \end{itemize}
\item
  (4)\(\iff\)(5)

  \begin{itemize}
  \item
    When \(r\pr{x}>0\), \(r\pr{x}\) is non-increasing in \(x\) if and only if 
    \(\ln r\pr{x} = \ln u'\pr{x} - \ln v'\pr{x}\) is non-increasing in
    \(x\).
  \item
    When \(u\) and \(v\) are twice differentiable,

    \[\pr{\ln r\pr{x}}' = \frac{u''\pr{x}}{u'\pr{x}} - \frac{v''\pr{x}}{v'\pr{x}}\]
  \item
    It follows that \(r\pr{x} = \frac{u'\pr{x}}{v'\pr{x}}\) is
    non-increasing in \(x\) if and only if \(A\pr{x,u}\ge A\pr{x,v}\).
  \end{itemize}
\end{itemize}
}

\ex{
Consider an investor who must decide how much of their
initial wealth \(w\) to put into a risky asset. The risky asset can have
any of the positive or negative rates of return \(r_i\) with
probabilities \(p_i\), \(i = 1, 2,\cdots, n\). Suppose the investor is
an expected utility maximizer and their utility for \(x\) amount of
money for sure can be represented by a twice continuously
differentiable, strictly increasing and strictly concave utility
function \(u\pr{x}\). Let \(a^*\) be the investor's optimal amount of
money to put in the risky asset. Suppose \(\sum_{i  =1}^n p_ir_i>0\).

\begin{enumerate}

\item
  Give a sufficient condition for \(a^*<w\).
\item
  Suppose \(a^*<w\) and that \(A\pr{x,u}\) is strictly decreasing in
  \(x\), how would \(a^*\pr{w}\) change with \(w\)?
\end{enumerate}
}


\hypertarget{comparison-of-risky-prospects}{%
\section{Comparison of Risky
Prospects}\label{comparison-of-risky-prospects}}

The earlier proposition shows we can use the coefficient of absolute
risk aversion \(A\pr{x,u}\) to compare the risk attitudes of two
individuals, as well as the risk attitudes of the same individual at two
wealth levels. But how about the comparisons of two lotteries? In
particular, when can we unambiguously say that \(F\pref G\) for any
(risk-averse) decision-maker?

When comparing monetary lotteries, an agent would naturally consider expected payoffs and variations in payoffs. Intuitively, \(F\pref G\) unambiguously in one of the following two
cases:

\begin{enumerate}
\item
  \(F\) yields a higher return than \(G\) with a higher probability.
\item
  \(F\) and \(G\) have the same expected return, but there is less
  variation in \(F\).
\end{enumerate}

Roughly speaking, we introduce first-order stochastic dominance for the first possibility, and second-order stochastic dominance for the second.

\hypertarget{first-order-stochastic-dominance}{%
\subsection{First-Order Stochastic
Dominance}\label{first-order-stochastic-dominance}}

\propp{First-Order Stochastic Dominance}{
The
following two conditions are equivalent:
\begin{enumerate}
\item
  For any non-decreasing \(u:\RR\to \RR\),
  \[\int_Xu\pr{t}\text dF\pr{t}\ge \int_X u\pr{t}\text dG\pr{t}.\]
\item
  \(F\pr{x}\le G\pr{x}\) for almost every $x$.
\end{enumerate}
In this case, \(F\) \textit{first-order stochastically dominates}
\(G\).
}{
Assume that \(u\) is continuously differentiable,
and \(u'\pr{\cdot}>0\). Let \(X = \br{\underline x,\bar x}\).
\[\ba
& \int_{\underline x}^{\bar x}u\pr{t}\text d F\pr{t}-\int_{\underline x}^{\bar x}u\pr{t}\text d G\pr{t}\\
=& \br{u\pr{t}F\pr{t}\bigg|_{\underline x}^{\bar x} - \int_{\underline x}^{\bar x} u'\pr{t}F\pr{t}\text d t}-\br{u\pr{t}G\pr{t}\bigg|_{\underline x}^{\bar x} - \int_{\underline x}^{\bar x} u'\pr{t}G\pr{t}\text d t}\\
= & \br{\pr{u\pr{\underline x}-u\pr{\bar x}}-\int_{\underline x}^{\bar x} u'\pr{t}F\pr{t}\text d t}-\br{\pr{u\pr{\underline x}-u\pr{\bar x}}-\int_{\underline x}^{\bar x} u'\pr{t}G\pr{t}\text d t} \\
= & \int_{\underline x}^{\bar x} u'\pr{t}\pr{G\pr{t}-F\pr{t}}\text dt
\ea\]

It follows that
\(\int_{\underline x}^{\bar x}u\pr{t}\text d F\pr{t}\ge \int_{\underline x}^{\bar x}u\pr{t}\text d G\pr{t}\)
if and only if \(G\pr{x}\ge F\pr{x}\) almost everywhere.
}

Notice that the first condition says that any decision-maker who prefers
more money would prefer lottery \(F\) over \(G\), regardless of the
agent's risk attitude, let alone the specific mathematical expression of utility function.

\hypertarget{second-order-stochastic-dominance}{%
\subsection{Second-Order Stochastic
Dominance}\label{second-order-stochastic-dominance}}

\defn{Mean-Preserving Spread}{
Let \(X,Y\) be two random
variables with distribution functions \(F\) and \(G\), respectively.
Then \(G\) is a \emph{mean-preserving spread} of \(F\) if
\(Y = X+\tilde \ve\), where \(\E{\tilde\ve|X} = 0\).
}

Intuitively, the two lotteries \(F\) and \(G\) yield the same expected
return, but \(G\) involves more risk.

\propp{Second-Order Stochastic Dominance}{
Suppose
\(\E{F} = \E{G}\). Then the following three conditions are equivalent:
\begin{enumerate}
\item
  For any (non-decreasing) concave \(u:\RR\to \RR\),
  \[\int_Xu\pr{t}\text dF\pr{t}\ge \int_Xu\pr{t}\text dG\pr{t}\]
\item
  For almost every \(x\in X\),
  \[\int_{-\infty}^x F\pr{t}\text dt\le \int_{-\infty}^x G\pr{t}\text dt\]
\item
  \(G\) is a mean-preserving spread of \(F\).
\end{enumerate}

In this case, \(F\) \textit{second-order stochastically
dominates} \(G\).
}{
For simplicity, suppose
\(X = \br{\underline x,\bar x}\) and \(u\) is twice continuously
differentiable.
\begin{itemize}
\item
  (1) \(\iff\) (2)
  \begin{itemize}
  \item
    By earlier calculation,
    \(\int_{\underline x}^{\bar x}u\pr{t}\text d F\pr{t}-\int_{\underline x}^{\bar x}u\pr{t}\text d G\pr{t}=\int_{\underline x}^{\bar x} u'\pr{t}\pr{G\pr{t}-F\pr{t}}\text dt\).
  \item
    Again by integration by parts and noticing \(\E{F} = \E{G}\), we
    have:
    \[\int_{\underline x}^{\bar x}u\pr{t}\text d F\pr{t}-\int_{\underline x}^{\bar x}u\pr{t}\text d G\pr{t}
    =\int_{\underline x}^{\bar x} u''\pr{x}\br{\int_{\underline x}^xG\pr{t}\text dt - \int_{\underline x}^x F\pr{t} \text dt} \text dx\]
  \item
    Because $u$ is concave, $u''\le 0$. It follows that
    \(\int_{\underline x}^{\bar x}u\pr{t}\text d F\pr{t}\ge \int_{\underline x}^{\bar x}u\pr{t}\text d G\pr{t}\)
    if and only if 
    \(\int_{\underline x}^xF\pr{t}\text dt \ge \int_{\underline x}^x G\pr{t} \text dt\)
    almost everywhere.
  \end{itemize}
\item
  (1) \(\iff\) (3)
  \begin{itemize}
  \item
    We will only show (3) \(\too\) (1).
  \item
    By the law of iterated expectation,
     \[\E[G]{u\pr Y} =  \E[F]{\E[G]{u\pr{Y}|X}}\]
  \item
    Since \(u\pr\cdot\) is concave, by Jensen's
    inequality,
    \[\E[G]{u\pr{Y}} = \E[F]{\E[G]{u\pr{Y}|X}} \le \E[F]{u\pr{\E[G]{Y|X}}} = \E[F]{u\pr{X}}\]
  \end{itemize}
\end{itemize}
}

\rmk{
The first condition says that any risk-averse decision-maker would prefer
  lottery \(F\) over \(G\). Additionally, when \(\E{F} = \E{G}\), we actually do not
  need \(u\pr{\cdot}\) to be non-decreasing. If \(\E{F}\ne \E{G}\), then the second and the third conditions are still
  equivalent, but we actually need \(u\pr{\cdot}\) to be non-decreasing
  and concave.
}

\hypertarget{likelihood-ratio-dominance}{%
\subsection{Likelihood Ratio
Dominance}\label{likelihood-ratio-dominance}}

\defn{Likelihood Ratio Dominance}{
Let \(F\) and \(G\) be two
distribution functions with common support \(\br{\underline x,\bar x}\).
Suppose the density functions exist and are given by \(f\) and \(g\),
respectively. \(F\) \textit{dominates} \(G\) \textit{in the likelihood ratio order}
if \(\frac{f\pr{x}}{g\pr{x}}\) is non-decreasing in \(x\).
}

Intuitively, when $F$ dominates $G$ in the likelihood ratio order, \(F\) puts higher probabilities on
higher returns compared with \(G\).

\propp{Likelihood Ratio Dominance Implies First-Order Stochastic Dominance}{
If \(F\)
dominates \(G\) in the likelihood ratio order, then \(F\) first-order
stochastically dominates \(G\).
}{
Since \(F\pr{\cdot }\) and \(G\pr{\cdot}\) are
both non-decreasing, in order to show that \(F\pr{x}\le G\pr{x}\), it
suffices to show that
\[\dfrac{F\pr{x}}{1-F\pr{x}}\le \dfrac{G\pr{x}}{1-G\pr{x}},\forall x\in \pr{\underline x,\bar x}.\]

We have
\[\ba
\dfrac{F\pr{x}}{1-F\pr{x}} & = \frac{\int_{\underline x}^x f\pr{t} \text dt}{\int_x^{\bar x}f\pr{t}\text dt}\\
& = \frac{\int_{\underline x}^x \frac{f\pr{t}}{g\pr{t}}\cdot g\pr{t} \text dt}{\int_x^{\bar x}\frac{f\pr{t}}{g\pr{t}}\cdot g\pr{t} \text dt}\\
\ea\]

Since \(\frac{f\pr{x}}{g\pr{x}}\) is non-decreasing,
\[\ba
\forall t\in \br{\underline x,x},\frac{f\pr{t}}{g\pr{t}}\le \frac{f\pr{x}}{g\pr{x}}\\
\forall t\in \br{x,\bar x},\frac{f\pr{t}}{g\pr{t}}\ge \frac{f\pr{x}}{g\pr{x}}\\
\ea\]
Therefore, we have
\[\ba
\dfrac{F\pr{x}}{1-F\pr{x}} 
& = \frac{\int_{\underline x}^x \frac{f\pr{t}}{g\pr{t}}\cdot g\pr{t} \text dt}{\int_x^{\bar x}\frac{f\pr{t}}{g\pr{t}}\cdot g\pr{t} \text dt}\le 
\frac{\frac{f\pr{x}}{g\pr{x}}\cdot\int_{\underline x}^x  g\pr{t} \text dt}{\frac{f\pr{x}}{g\pr{x}}\cdot \int_x^{\bar x}g\pr{t} \text dt} = 
\frac{\int_{\underline x}^x  g\pr{t} \text dt}{\int_x^{\bar x}g\pr{t} \text dt} = \frac{G\pr{x}}{1-G\pr{x}}
\ea\]
}

\hypertarget{comparative-statics-under-risk}{%
\section{Comparative Statics Under
Risk}\label{comparative-statics-under-risk}}

Our earlier discussions on comparative statics did not 
specifically take uncertainty into consideration. Whether and how would the comparative
statics results carry over to choices under uncertainty?

The first result is that, if a random variable is ``introduced" into a supermodular function, the expected utility function still preserves supermodularity.

\lemp{Expected Utility Function Preserves Supermodularity}{
Let
\(X\subset \RR^n\) be a sublattice and \(T\subset \RR\). Suppose
\(u:X\times T\to \RR\) is supermodular in \(\x\). Then for any
distribution function \(F\) on \(T\), the function \(U:X\to \RR\) defined
by
\[U\pr{\x} = \int_t u\pr{\x,t}\text dF\pr{t}\]
\hspace{2em} is supermodular in \(\x\).
}{
Take any \(\x,\x'\in X\). we have
\[\ba
U\pr{\x\meet \x'}+U\pr{\x\join \x'} & = \int_t\br{u\pr{\x\meet \x',t}+u\pr{\x\join \x',t}}\text d F\pr{t}\\
& \ge \int_t\br{{u\pr{\x,t}+u\pr{\x',t}}}\text d F\pr{t}\\
& = U\pr{\x}+U\pr{\x'}
\ea\]

where the inequality in the second line follows from the supermodularity of \(u\pr{\x,t}\)
in \(\x\). Consequently by definition of supermodularity, \(U\pr{\cdot}\) is supermodular in \(\x\).
}

The second result states that, first-order stochastic dominance preserves increasing differences.

\lemp{First-Order Stochastic Dominance Preserves Increasing Differences}{
Let \(X\subset \RR^n\)
be a sublattice and \(T,\Theta\subset \RR\). Suppose
\(u:X\times T\to \RR\) has increasing differences in \(\pr{\x,t}\), and
\(\brace{F_{\theta}}_{\theta\in \Theta}\) is a family of distribution
function on \(T\) such that \(F_{\theta}\ge_{FOSD} F_{\theta'}\) if
\(\theta>\theta'\). Then the function \(U:X\times\Theta\to \RR\) defined
by
\[U\pr{\x,\theta} = \int_t u\pr{\x,t}\text dF_{\theta}\pr{t}\]
\hspace{2em} has increasing differences in \(\pr{\x,\theta}\).
}{
Take any \(\x>\x'\) and \(\theta>\theta'\), we have
\[\ba
U\pr{\x,\theta}-U\pr{\x',\theta} & = \int_t \br{u\pr{\x,t}-u\pr{\x',t}}\text dF_{\theta}\pr{t}\\
& \ge \int_t \br{u\pr{\x,t}-u\pr{\x',t}}\text dF_{\theta'}\pr{t}\\
& = U\pr{\x,\theta'}-U\pr{\x',\theta'}
\ea\]

where the inequality holds because, \(\delta\pr{t}:= u\pr{\x,t}-u\pr{\x',t}\) is non-decreasing in \(t\)
by increasing differences of \(u\pr{\cdot,\cdot}\) in \(\pr{\x,t}\) and 
\(F_{\theta}\ge_{FOSD} F_{\theta'}\).
Consequently, \(U\pr{x,\theta}\) has increasing differences in
\(\pr{\x,\theta}\).
}

\rmk{
Intuitively, \(\theta\) can be interpreted as a signal of \(t\), indicating the distribution of $t$.
}

\propp{Comparative Statics Under Risk}{
Let
\(X\subset \RR^n\) be a sublattice and \(T,\Theta\subset \RR\). Suppose
\(u:X\times T\to \RR\) is supermodular in \(\x\) and has increasing
differences in \(\pr{\x,t}\). Further suppose
\(\brace{F_{\theta}}_{\theta\in \Theta}\) is a family of distribution
function on \(T\) such that \(F_{\theta}\ge_{FOSD} F_{\theta'}\) if
\(\theta>\theta'\). Then for the function \(U:X\times\Theta\to \RR\)
defined by
\[U\pr{\x,\theta} = \int_t u\pr{\x,t}\text dF_{\theta}\pr{t},\]
\hspace{2em} the set of maximizers, \(\arg\max_{\x\in X} U\pr{\x,\theta}\), is
non-decreasing in \(\theta\) in the strong set order.
}{
\begin{itemize}
\item
  Fixing any \(\theta\), by the lemma of expected utility function preserving supermodularity, \(U\pr{\x,\theta}\) is
  supermodular in \(\x\).
\item
  By the lemma of first-order stochastic dominance preserving increasing differences, \(U\pr{\x,\theta}\) has increasing differences in
  \(\pr{\x,\theta}\).
\item
  Since \(X\times \Theta\) forms a product set, \(U\pr{\x,\theta}\) is
  supermodular in \(\pr{\x,\theta}\).
\item
  By the multivariate Topkis' theorem, the set of
  maximizers, \(\arg\max_{\x\in X} U\pr{\x,\theta}\), is non-decreasing in
  \(\theta\) in the strong set order.
\end{itemize}
}

\ex{
Suppose a monopolist faces an uncertain demand and must
make a production decision prior to learning the realized demand for its
product. The monopolist does learn some information about demand prior
to choosing its output. Formally, suppose the inverse demand function is 
\(p\pr{q,t} = \hat p\pr{q}+t\), where \(\hat p\pr{q}\) is the estimated
inverse demand, and \(t\) is a random noise. The ex-post profit of the
firm is therefore
\[\pi\pr{q,t} = p\pr{q,t}\cdot q-c\pr{q}\]
\hspace{2em} with \(c\pr q\) being the monopolist's cost function. The
monopolist observes a signal \(\theta\) that is informative about the
parameter \(t\). Specifically, suppose the distribution of \(t\)
conditional on \(\theta\) is \(F_{\theta}\pr{\cdot}\), and
\(F_{\theta}\ge_{FOSD} F_{\theta'}\) for \(\theta>\theta'\). Determine
how would the monopolist's optimal output
\(q^*\pr{\theta}\) change with the observed signal \(\theta\).
}

\sol{
The monopolist solves the following maximization
problem:
\[\max_{q\ge 0} \Pi \pr{q,\theta} = \int_t \pi\pr{q,t}\text d F_{\theta}\pr t\]

Notice that \(\frac{\partial \pi\pr{q,t}}{\partial t} = q\), increasing in \(q\). It follows that \(\pi\pr{q,t}\) has increasing differences in
\(\pr{q,t}\). Since \(F_{\theta}\ge_{FOSD} F_{\theta'}\) for
\(\theta>\theta'\), \(\Pi\pr{q,\theta}\) has increasing differences in
\(\pr{q,\theta}\) and \(q^*\pr{\theta}\) is non-decreasing in \(\theta\)
in the strong set order.
}

\ex{
Consider an agent with an initial wealth \(w\), who
faces a random loss \(\tilde l\in \br{0,w}\). To counter the potential
loss, the agent may purchase any fraction of insurance
\(y\in \br{0,1}\). Specifically, if the agent purchases \(y\) unit of
insurance, then they pay a total price of \(yp\) upfront, and get paid
\(yl\) if they incur a loss of \(l\). Suppose the agent is an expected
utility maximizer and has a utility function \(u\pr{x}\) for $x$ 
amount of money for sure. Further suppose \(u\pr{\cdot}\) is twice
continuously differentiable with \(u'\pr{\cdot}>0\) and
\(u''\pr{\cdot}<0\).

\begin{enumerate}
\item
  First suppose \(p\le \E{\tilde l}\). Show the optimal amount of
  insurance \(y^*=  1\).
\item
  Next suppose \(p>\E{\tilde l}\). Show the optimal amount of insurance
  \(y^*<1\).
\item
  Let \(A\pr{x,u}\) be the coefficient of absolute risk aversion. Write
  down its expression and show that if \(A\pr{x,u} = c_0\) (a constant),
  then the optimal amount of insurance \(y^*\) is independent of \(w\).
\item
  Now suppose \(A\pr{x,u}\) strictly decreases with \(x\). Show the
  optimal amount of insurance \(y^*\) is non-increasing in the
  agent's initial wealth \(w\).
\end{enumerate}
}
\sol{
\begin{enumerate}
    \item Since the agent is an expected utility maximizer, their choice problem
is given by
\[\max_{y\in \br{0,1}} U\pr y = \E{u\pr{w-yp-\pr{1-y}l}}\]

Simple calculation shows
\[\ba
U'\pr y & = \E{\pr{l-p}\cdot u'\pr{w-l+\pr{l-p}y}}\\
U''\pr y & = \E{\pr{l-p}^2\cdot u''\pr{w-l+\pr{l-p}y}}\le 0
\ea\]

It follows that
\(U'\pr{y}\ge U'\pr1 = u'\pr{w-p}\pr{\E{\tilde l}-p}\ge 0\), with strict
inequality as long as 
\(p<\E{\tilde l}\). Consequently, \(y^* = 1\).

\item  When \(p>\E{\tilde l}\), \(U'\pr1 = u'\pr{w-p}\pr{\E{\tilde l}-p}<0\), and \(U''\pr y\le 0\), we have \(y^*<1\).

\item By definition, \(A\pr{x,u} = -\dfrac{u''\pr x}{u'\pr x}\).

Take any \(w_1<w_2\), and let
\[\ba
v_1\pr x & = u\pr{w_1+x}\\
v_2\pr x & = u\pr{w_2+x}
\ea\]
Then at initial wealth \(w_i\pr{i = 1,2}\), the agent's
maximization problem is given by
\[\max_{y\in \br{0,1}}V_i\pr y = \E{v_i\pr{-l+\pr{l-p}y}}\]
From the previous two questions, it suffices to focus on the case of
\(p>\E{\tilde l}\).

By earlier characterization, \(A\pr{x,u} = c_0\) implies
\(\dfrac{v_1'\pr x}{v_2'\pr x} = c_1\), which is also a constant.

Simple calculation shows
\[\ba
V_i'\pr y & = \E{\pr{l-p}v_i'\pr{-l+\pr{l-p}y}}\\
V_i''\pr y & = \E{\pr{l-p}^2v_i''\pr{-l+\pr{l-p}y}}\le 0
\ea\]
If \(V_1'\pr 0\le 0\), then \(V_2'\pr{0}\le 0\), so
\(y_1^* = y_2^* = 0\). Otherwise, if \(V_1'\pr{y^*} = 0\), then
\(V'_2\pr{y^*} = 0\); by $V_i''(\cdot)< 0$ for $i = 1,2$, $V_i'(y)$ strictly decreases with $y$ and achieves maximum at $y$, so \(y_1^* = y_2^* = y^*\).

\item Suppose by contradiction that \(y_2^*>y_1^*\ge 0\). By optimality of $y_2^*$, 
\[\int_0^w v_2'\pr{-l+\pr{l-p}y_2^*}\pr{l-p}\text dF\pr{l} = 0.\]
By our earlier
characterization, \(A\pr{x,u}\) strictly decreases with \(x\) implies
\(\dfrac{v'_1\pr x}{v_2'\pr x}\) is non-increasing in \(x\).

Consequently, we have
\[\ba
V_1'\pr{y_2^*} & = \int_0^p v_1'\pr{-l+\pr{l-p}y_2^*}\pr{l-p}\text dF\pr{l}\\
& \ \ \ \ +\int_p^w v_1'\pr{-l+\pr{l-p}y_2^*}\pr{l-p}\text dF\pr{l}\\
& = \int_0^p \frac{v_1'\pr{-l+\pr{l-p}y_2^*}}{v_2'\pr{-l+\pr{l-p}y_2^*}}v_2'\pr{-l+\pr{l-p}y_2^*}\pr{l-p}\text dF\pr{l}\\
& \ \ \ \ +\int_p^w \frac{v_1'\pr{-l+\pr{l-p}y_2^*}}{v_2'\pr{-l+\pr{l-p}y_2^*}}v_2'\pr{-l+\pr{l-p}y_2^*}\pr{l-p}\text dF\pr{l}\\
& \ge \frac{v_1'\pr{-p}}{v_2'\pr{-p}}\int_0^w v_2'\pr{-l+\pr{l-p}y_2^*}\pr{l-p}\text dF\pr{l}\\
& = 0
\ea\]
Given $V_i''(\cdot)<0$ and our assumption of $y_2^*>y_1^*\ge 0$, we have
\(V_1'\pr{y_1^*}> V_1'\pr{y_2^*}\ge 0\), which is a contradiction.
\end{enumerate}
}

